\subsection{Railroad Inventory}
\subsubsection{aarcodes}\label{aarcode}
An \gls{aar} code refers to a classification code used by the \gls{aar} to identify different types of equipment and rollingstock used in the North American railway industry. The AAR codes are a standardized way of categorizing and labeling various railway vehicles, making it easier to identify and track equipment across different railways and organizations.
\begin{longtable}{|L{3cm}|L{1.7cm}|L{9.1cm}|L{1.0cm}|}
	\caption{\label{aar-table}AarCode Collection Fields Table}\\
    \hline
	\textbf{Field} & \textbf{Type} & \textbf{Description} & \textbf{Req'd} \\
	\hline
	\endfirsthead
	\multicolumn{4}{c}%
	{\tablename\ \thetable\ -- \textit{Continued from previous page}} \\
	\hline
	\textbf{Field} & \textbf{Type} & \textbf{Description} & \textbf{Req'd} \\
	\hline
	\endhead
	\hline \multicolumn{4}{r}{\textit{Continued on next page}} \\
	\endfoot
	\hline
	\endlastfoot
        aarCode & String & A classification code used by the \gls{aar} to identify different types of rollingstock & Yes \\ \hline
        rollingstockType & String & A category of rollingstock having common characteristics, such as; Auxiliary, Freight, Gondola, Hopper, Locomotive, Non Revenue, and Passenger. & No\\ \hline
        description & String & A written representation that aims to make vivid the characteristics of the AAR Code. & No \\ \hline
\end{longtable}
\subsubsection{dccs}
\gls{dcc} is a system used in model railroading to control and operate model trains and their accessories. It is a standardized method of transmitting control signals digitally to locomotives and other rollingstock. \gls{dcc} allows for more precise and independent control of multiple trains on the same track.
\begin{longtable} { |L{3cm}|L{1.7cm}|L{9.1cm}|L{1.0cm}| }
	\caption{\label{dcc-table}Dccs Collection Fields Table}\\
    \hline
	\textbf{Field} & \textbf{Type} & \textbf{Description} & \textbf{Req'd} \\
	\hline
	\endfirsthead
	\multicolumn{4}{c}%
	{\tablename\ \thetable\ -- \textit{Continued from previous page}} \\
	\hline
	\textbf{Field} & \textbf{Type} & \textbf{Description} & \textbf{Req'd} \\
	\hline
	\endhead
	\hline \multicolumn{4}{r}{\textit{Continued on next page}} \\
	\endfoot
	\hline
	\endlastfoot
        locomotiveID & ObjecctId &  The identity object of the dccs document assigned by MongoDB at the time the rollingstock (locomotive) document is created.& Yes \\ \hline
        mfg & String & The company that made the \gls{dcc} decoder installed in a rollingstock, typically a locomotive. & No\\ \hline
        family & String & A group of DCC decoders produced by a specific manufacturer that share common features, specifications, and compatibility. & No \\ \hline
        model &  String & A \gls{dcc} decoder offering specific features or specifications within the same product family& No \\ \hline
        address & String & A unique identifier assigned to each \gls{dcc} decoder, which is a numerical or digital label that distinguishes one rollingstock from another in the \gls{dcc} system. \gls{dcc} addresses range from 1 to 9999. & No \\ \hline  
\end{longtable}
\subsubsection{images}\label{images}
A set of photographs or illustrations of the rollingstock and structures of the model railroad present in \gls{rails}. These images may capture different designs, configurations, and features of rollingstock, providing a visual representation of the diversity within the railway transportation system.
\begin{longtable}{|L{3cm}|L{1.7cm}|L{9.1cm}|L{1.0cm}|}
	\caption{\label{image-table}Images Collection Fields Table}\\
	\hline
    \textbf{Field} & \textbf{Type} & \textbf{Description} & \textbf{Req'd} \\
	\hline
	\endfirsthead
	\multicolumn{4}{c}%
	{\tablename\ \thetable\ -- \textit{Continued from previous page}} \\
	\hline
	\textbf{Field} & \textbf{Type} & \textbf{Description} & \textbf{Req'd} \\
	\hline
	\endhead
	\hline \multicolumn{4}{r}{\textit{Continued on next page}} \\
	\endfoot
	\hline
	\endlastfoot
        fileName & String & A unique identifier assigned to a computer file in order to distinguish it from other files within a file system. & Yes \\ \hline
        title & String & The name or label given to the image to identify and distinguish it from others. & No\\ \hline
        notes & String & A brief record of events or observations about the image. & No \\ \hline
        category & String & A classification or group of images that share common characteristics, features, or attributes. & No \\ \hline
\end{longtable}
\subsubsection{industries}
Sectors of economic activity that involve the production of goods or the provision of services. These sectors are characterized by similar business activities and are often grouped together based on their common features. Railways service a wide range of industries, facilitating the transportation of goods and passengers. The industries collection is used to identify the industries that are present in the model railroad system.
\begin{longtable}{|L{3cm}|L{1.7cm}|L{9.1cm}|L{1.0cm}|}
	\caption{\label{industry-table}Industries Collection Fields Table}\\
	\hline
	\textbf{Field} & \textbf{Type} & \textbf{Description} & \textbf{Req'd} \\
	\hline
	\endfirsthead
	\multicolumn{4}{c}%
	{\tablename\ \thetable\ -- \textit{Continued from previous page}} \\
	\hline
	\textbf{Field} & \textbf{Type} & \textbf{Description} & \textbf{Req'd} \\
	\hline
	\endhead
	\hline \multicolumn{4}{r}{\textit{Continued on next page}} \\
	\endfoot
	\hline
	\endlastfoot
        shortName & String & A concise or abbreviated version of a longer name. It could be a nickname, an acronym, or a shortened form that is used for convenience or brevity. & Yes \\ \hline
        longName & String &  A word or set of words by which an industry is identified and distinguished from others, legal name. & No\\ \hline
        industryType & String & The various categories or sectors of economic activity in which businesses and organizations engage. & No \\ \hline
        industryLocation & String & A specific place or position in space, often designated by its geographical coordinates or described in relation to other landmarks. & No \\ \hline
\end{longtable}
\subsubsection{rollingstocks}
rollingstock refers to the vehicles that move on a railway track. These vehicles are an essential part of any railway system and include various types used for transporting goods and passengers. Rollingstock can be broadly categorized into two main types: passenger rollingstock and freight rollingstock. Locomotives are considered a type of rollingstock in the context of railway transportation.
\begin{longtable}{|L{3cm}|L{1.7cm}|L{9.1cm}|L{1.0cm}|}
	\caption{\label{rollingstock-table}Rollingstock Collection Fields Table}\\
    \hline
    \textbf{Field} & \textbf{Type} & \textbf{Description} & \textbf{Req'd} \\
	\hline
	\endfirsthead
	\multicolumn{4}{c}%
	{\tablename\ \thetable\ -- \textit{Continued from previous page}} \\
	\hline
	\textbf{Field} & \textbf{Type} & \textbf{Description} & \textbf{Req'd} \\
	\hline
	\endhead
	\hline \multicolumn{4}{r}{\textit{Continued on next page}} \\
	\endfoot
	\hline
    \multicolumn{4}{|p{14.8cm}|}{\small \textit{Note:} The numbers in parentheses in the Req'd column are the default values.} \\
	\hline
	\endlastfoot
	roadName & String & The road name also reffered to as reporting marks are a set of letters used to identify the owner or operator of a railcar or locomotive. & Yes \\ \hline
	roadNumber & String & The road number is an identifier for a railcar or a locomotive and is numerical. & Yes \\ \hline
	color & String & The exterior paint or livery applied to the rollingstock. The color of rollingstock serves several purposes, including identification, branding, and visibility. & No \\ \hline
	aarCode & String & See paragraph \ref{aarcode}  & No \\ \hline
	description & String &  The informative account that provides information about the characteristics, features, qualities, or attributes of the rollingstock. & No \\ \hline
	numberBlt & Long & The total quantity or count of a particular rollingstock. & No (0)\\ \hline
	inSvcDate & Date & The specific date on which a peice of rollingstock is officially put into operation or use. It marks the beginning of its active service life, during which it is expected to perform its intended functions. & No \\ \hline
	insideLength & String & For passenger cars it is measured from one end of the passenger compartment to the other. It defines the available space for seating and other passenger-related amenities. For freight cars it is measured from the inside faces of the end walls, excluding any bulkheads or projections that might reduce the usable length. It provides an indication of the space available for loading and transporting goods. & No \\ \hline
	insideHeight & String &  & No \\ \hline
	insideWidth & String &  & No \\ \hline
	loadTypes & String &  & No \\ \hline
	capacity & Long & Capacity refers to the maximum amount of cargo, typically measured in weight or volume, that it can safely and efficiently transport. & No (0)\\ \hline
	bldr & String &  & No \\ \hline
	bltDate & Date &  & No \\ \hline
	notes & String & A brief record of events or observations about the rollingstock. & No \\ \hline
	ltWeight & Long & The weight of the car itself without any cargo or load. & No (0)\\ \hline
	loadLimit & Long & The maximum weight that a freight car can carry, including both the weight of the car itself (tare weight or lightweight) and the weight of the cargo (payload). & No (0)\\ \hline
	lastMaintDate & Date &  & No \\ \hline
	locationNow & String &  & No \\ \hline
	homeLocation & String &  & No \\ \hline
	rsStatus & String & The current condition, availability, or operational state of the rollingstock used on a railway. & No \\ \hline
    numAxles & Long &  The total count of transverse metal rods or central shafts that span beneath the rollingstock to support its weight and connect opposite wheel assemblies. & No (4)\\ \hline
	imageID & String &  The file name of the image; See paragraph \ref{images} & No \\ \hline
	modelWeight & Double & The mass of the model rollingstock in either grams or ounces. & No (0)\\ \hline
	modelLength & Double & The distance measure between the middle of the front coupler knuckle to the middle of the rear coupler knuckle of the model rollingstock. & No (0)\\ \hline
	rfid & String & The 10 character identifier emitted by the RFID tag attached to the model rollingstock. & No \\ \hline
    rfidLocation & Double & The distance the RFID tag is attached to the rollingstock measured from the A end. & No (0)\\ \hline
    \end{longtable}
\subsubsection{structures}
\begin{table}[H]
    \begin{tabular}{|L{3cm}|L{1.7cm}|L{8.8cm}|L{1.3cm}|}
    \hline
        \textbf{Field} & \textbf{Type} & \textbf{Description} & \textbf{Req'd} \\ \hline
	title & String &  & Yes \\ \hline
	structureUse & String &  & Yes \\ \hline
	description & String &  & No \\ \hline
	owner & String &  & No \\ \hline
	location & String &  & No \\ \hline
	construction & String &  & No \\ \hline
	builtDate & String &  & No \\ \hline
	size & String &  & No \\ \hline
	image & String &  & No \\ \hline
	isIndustrial & Boolean &  & No \\ \hline
	rawMaterials & String &  & No \\ \hline
	rMCapacity & String &  & No \\ \hline
	conRate & String &  & No \\ \hline
	priority & String &  & No \\ \hline
	aarCodeIn & String &  & No \\ \hline
	product & String &  & No \\ \hline
	productCap & String &  & No \\ \hline
	prodRate & String &  & No \\ \hline
	aarCodeOut & String &  & No \\ \hline
	unloadDuration & String &  & No \\ \hline
	loadDuration & String &  & No \\ \hline
	sidingCap & String &  & No \\ \hline
	notes & String &  & No \\ \hline
    \end{tabular}
    \caption{\label{structure-table}Structures Collection Fields Table}
    \end{table}

